\documentclass[12pt]{article}
\usepackage[paperwidth=5.5in,paperheight=3.2in,margin=0.4in]{geometry}
\usepackage{amsmath}
\usepackage{tikz}
\pagestyle{empty}
\begin{document}
\noindent
Consider $\sum_{k=1}^n k = n(n+1)/2$ for now. A text-style operator such as
$\sum$ or $\int_0^1 f(x)\,dx$ is centered on the math axis, so it rises above
the cap height and dips below the baseline, just like the descenders in
``Copy gym''. Fractions like $\frac{a}{b}$ shrink to script size inline, while
a tall delimiter $\left(\frac{1}{2}\right)$ simply overflows the line and
\TeX\ pushes the following baseline down to make room. Plain symbols such as
$x^2 + y_i$ and $\sqrt{x+1}$ stay put on the baseline.

\medskip
\noindent
Tight box around the line: \fbox{Consider $\sum_{k=1}^n k = n(n+1)/2$ for now.}

\medskip
\noindent
Strut-sized line box drawn: \tikz[baseline]{\draw[gray,thin] (0,-0.3em) rectangle (17.2em,0.7em); \node[anchor=base west] at (0,0) {Consider $\sum_{k=1}^n k = n(n+1)/2$ for now.};}
\end{document}
